% ==================================================================
% id:       cond-027
% topic:    הסתברות מותנית ואי-תלות
% subtopic: מאורעות בלתי תלויים
% solution: full
% ==================================================================
במעבדה לביולוגיה בודקים איך דשן משפיע על נביטה של זרעים. $20\%$ מהזרעים נשתלים עם דשן מועשר, $30\%$ עם דשן רגיל, ו־$50\%$ בלי דשן. ההסתברות שזרע נובט היא $0.35$. מבין הזרעים שנשתלו בלי דשן, $80\%$ לא נבטו, ומבין הזרעים שנשתלו עם דשן רגיל, $75\%$ לא נבטו.
\begin{enumerate}
\item מה ההסתברות שזרע שנשתל עם דשן מועשר נובט?
\hint{כתבו את נוסחת ההסתברות השלמה עבור ״הזרע נבט״, כשהנעלם הוא ההסתברות המבוקשת.}
\begin{solbox}
נסמן $D_2$ = ״דשן מועשר״, $D_1$ = ״דשן רגיל״, $D_0$ = ״בלי דשן״ (חלוקה של $\Omega$), ו־$N$ = ״הזרע נבט״. נתון $\PP(N)=0.35$, $\PP(N\mid D_0)=1-0.8=0.2$ ו־$\PP(N\mid D_1)=1-0.75=0.25$. נסמן $x=\PP(N\mid D_2)$. לפי נוסחת ההסתברות השלמה
\[
\PP(N)=\PP(N\mid D_2)\,\PP(D_2)+\PP(N\mid D_1)\,\PP(D_1)+\PP(N\mid D_0)\,\PP(D_0),
\]
כלומר
\[
0.35=0.2x+0.3\cdot0.25+0.5\cdot0.2=0.2x+0.175,
\]
ולכן $0.2x=0.175$ ו־$x=\ans{0.875}$.
\end{solbox}
\item בוחרים זרע באקראי, ומתברר שהוא לא נבט. מה ההסתברות שהוא נשתל בלי דשן?
\begin{solbox}
$\PP(N^C)=1-0.35=0.65$, ולפי נוסחת בייס
\[
\PP(D_0\mid N^C)=\frac{\PP(N^C\mid D_0)\,\PP(D_0)}{\PP(N^C)}=\frac{0.8\cdot0.5}{0.65}=\frac{0.4}{0.65}=\ans{\frac8{13}}\approx0.615.
\]
\end{solbox}
\item נסמן $A$ = ״הזרע נשתל בלי דשן״ ו־$B$ = ״הזרע לא נבט״. האם $A$ ו־$B$ בלתי תלויים? האם הם זרים?
\begin{solbox}
לפי כלל הכפל $\PP(A\cap B)=\PP(B\mid A)\,\PP(A)=0.8\cdot0.5=0.4$, ואילו $\PP(A)\,\PP(B)=0.5\cdot0.65=0.325$. אלה שונים, ולכן $A$ ו־$B$ \(\ans{\text{תלויים}}\). הם גם \(\ans{\text{אינם זרים}}\), כי $\PP(A\cap B)=0.4>0$.
\end{solbox}
\end{enumerate}
